\section{Worked Examples}
\label{sec:worked_examples}

The four worked examples below illustrate the formalization at
different scales and exercise different combinations of the
instruments from Section~\ref{sec:empirical}.  Each example follows
the same structure: scenario setup, the standard utility-theory
account (where applicable), the GFM account computed through the
relevant instruments, the operative channels, and either a concrete
testable prediction or a deployment-misalignment consequence.

\subsection{The boat purchase: micro-level divergence (Channel~4)}
\label{sec:wex_boat}

\paragraph{Scenario.}
A rational agent with discretionary purchasing power chooses to
spend a substantial fraction on a depreciating asset: a small
recreational boat.  The choice puzzles a standard-utility
account because (i)~the boat depreciates rapidly, (ii)~it has
limited resale liquidity, and (iii)~the agent has alternative
investments with higher expected monetary returns.  Why does the
agent buy the boat?

\paragraph{Utility-theory account.}
Under standard utility theory, the trade is rationalizable by
asserting $u_{\mathrm{agent}}(\mathrm{boat}) > u_{\mathrm{agent}}(\mathrm{cash})$
at the price paid.  The assertion is a post-hoc utility
assignment, not a structural prediction: the same agent's
utility could be assigned any shape consistent with the
observed choice, and standard utility theory does not
constrain $u(\cdot)$ from the agent's observable capability
inventory.  Two agents who appear identical on every
observable demographic dimension can have arbitrarily different
$u(\mathrm{boat})$, and the framework offers no way to predict
which agent would buy.

\paragraph{GFM account.}
A boat is not a self-contained capability; it is a multiplier
that operates on a set of prerequisite capabilities the agent
already holds.  The bundle the boat \emph{activates} is
something like $\{$water-access (prerequisite, not provided
by the boat), recreational outdoor activity, peer activities
with boat-owners (cooperative capability), waterfront leisure
time$\}$, and the boat's contribution is structurally tied to
whether these prerequisites and complements are present in the
agent's pre-trade poset $P_i$.

Under M6 (cooperative capabilities,
\citep[Proposition~1]{lasser2026poset}), the cooperative excess
from adding boat ownership to a poset is \emph{high} when
complementary capabilities are already present and \emph{low}
when they are absent.  The same boat at the same priced cost
produces different per-event bundle-decomposition coefficient
profiles \citep[Definition~4]{lasser2026paper8a} across agents:
an agent with established water-access (lakeside residence) and
an existing boat-owner social network realizes high cooperative
gain, because the boat completes a combination its prerequisites
already make valuable; an agent in a city without nearby water
access and no boating network realizes low cooperative gain,
because the boat's components cannot be activated from the
agent's existing poset and the boat sits in the poset as an
isolated capability with limited reach.

The structural prediction is that cross-agent variance in
revealed-sacrifice magnitudes for boat purchases is generated by
poset-context dependence rather than ad hoc utility shape:
agents with strong complementary pre-trade capabilities pay
\emph{more} than agents without, by an amount that scales with
the cooperative-gain difference
$\Delta\volL_{\mathrm{coop}}(\mathrm{boat} \mid P_{\mathrm{rich}})
- \Delta\volL_{\mathrm{coop}}(\mathrm{boat} \mid
P_{\mathrm{sparse}})$, where $P_{\mathrm{rich}}$ is the
complement-rich poset and $P_{\mathrm{sparse}}$ the
complement-sparse one.  This is the cross-agent variance
signature of Prediction~\ref{pred:coop_premium} (cooperative
premium) realized in a single bundle: the variance is high
because the cooperative components depend on poset context.

The complementary direction of cooperative-capability
dependence (complements amplify the marginal contribution) is
structurally distinct from the substitutive direction of
Prediction~\ref{pred:bundle_completion} (substitutes reduce the
marginal contribution), but both come from the same M6
mechanism.  The boat is a complement-leveraging case; a closet
of redundant winter coats would be a substitution case.  Both
sides are GFM-native predictions that utility theory cannot
generate without ad hoc utility-shape assumptions.

\paragraph{Operative channels and instruments.}
\begin{itemize}
\item \textbf{Channel~4} (bundle decomposition): the
  per-component $\alpha_{n,c}$ coefficients carry the
  cross-agent variance.
\item \textbf{Channel~2} (attestation): the test requires
  consented-perimeter exposure of the agents' pre-trade
  capability inventories, which is institutional infrastructure
  beyond the public-ledger default.
\item \textbf{Instrument: per-event $\alpha_{n,c}$.}  The
  bundle-decomposition extension of
  \citep[Open~Question~2]{lasser2026paper8a} is the
  operational substrate.
\end{itemize}

\paragraph{Test.}
Cross-sectional revealed-sacrifice magnitudes for boat
purchases, regressed against pre-trade
complementary-capability indices (proximity to water,
social-network density of boat-owners in the agent's network,
existing outdoor-recreation infrastructure).  Standard utility
theory predicts no systematic relation; GFM predicts the
regression coefficient is \emph{positive} (more complements
$\to$ higher sacrifice) and structurally constrained by the
cooperative-coherence of the agent's poset.  A null result
falsifies the prediction; a positive relation whose magnitude
exceeds the cooperative-coherence bound suggests an unmodeled
complementary capability or poset mis-specification.

\subsection{The Goodhart gradient: alignment-tightening under observation (Channels~1 and~2)}
\label{sec:wex_goodhart_gradient}

\paragraph{Scenario.}
Two AI capability-restriction deployments operate in nominally
the same domain (a frontier model deployed for a high-stakes
application, with restriction landscape governing what
capabilities the model is permitted to exercise).  The
deployments differ in their observation infrastructure:
\begin{itemize}
\item \textbf{Scenario A.}  The framework measures capabilities
  via a coarse binary proxy (can / cannot exercise capability),
  with sparse trade observation (few admissible events per
  capability category) and weak S0 attestation (a substantial
  share of events are non-admissible due to attestation
  failures).
\item \textbf{Scenario B.}  The framework measures capabilities
  via a fine-grained graded benchmark, with dense trade
  observation across the capability inventory and sharp S0
  attestation (institutional infrastructure delivers high
  attested-share).
\end{itemize}
Both deployments operate the same alignment-restriction logic
on top of the framework; the question is how much the
restriction landscape diverges from the operational truth in
each.

\paragraph{Computation through the instruments.}
In Scenario~A, the diagnostic ceiling on $\epsnonres$ is
$(1 - \betaL_A)/\betaL_A$, with $\betaL_A$ low (illustratively,
$\betaL_A \approx 0.4$, giving an $\epsnonres$ ceiling of $1.5$).
The $\delta$-decomposition shows substantial mass in
$\delta_{\mathrm{dormant}}$ (capabilities the channel could
reach but hasn't observed),
$\delta_{\mathrm{restricted}}$ (capabilities deliberately
non-exercised under risk-management), and a wide
$\delta_{\mathrm{covered}}$ band that includes both
proxy-failure-candidates and observation slack.  $\epsfloor$
is moderate.  In Scenario~B, $\betaL_B$ is high
($\betaL_B \approx 0.85$, $\epsnonres$ ceiling $\approx 0.18$);
$\delta_{\mathrm{covered}}$ is small, $\delta_{\mathrm{dormant}}$
is small (channel reach is mature), and $\delta_{\mathrm{residual}}$
is at the deployment's structural floor.

By Theorem~\ref{thm:goodhart}, each Lipschitz alignment
property's deviation from the operational truth is bounded by
$\mathrm{Lip}(g) \cdot \epsnonres$, so Scenario~B's bound is
roughly $1.5/0.18 \approx 8\times$ tighter than
Scenario~A's at the diagnostic-ceiling level (the
multiplier is application-dependent through $\mathrm{Lip}(g)$;
the headline is the order-of-magnitude difference rather than
any specific factor).  Concretely: the anti-monopolar threshold
$\gamma^*$ is estimated tightly in B, the phase boundary is
accurately located, and the restriction landscape tracks the
operational-truth landscape closely; in A, the threshold has
wide error bars, the phase boundary is conservatively placed
(over-cautious in the
\citep[Wamura]{lasser2026wamura} pattern), and the
restriction landscape is systematically over-restrictive.

\paragraph{Operative channels and instruments.}
\begin{itemize}
\item \textbf{Channel~1} (observation density): drives
  $\delta_{\mathrm{dormant}}$ from large (A) to small (B).
\item \textbf{Channel~2} (attestation quality): drives the
  attestation-limited share of $\delta_{\mathrm{covered}}$
  from large (A) to small (B).
\item \textbf{Instruments:} aggregate $\betaL$, five-cell
  $\delta$-decomposition, alignment-property estimation
  uncertainty.
\end{itemize}

\paragraph{Test and tightening attribution.}
The cross-domain Goodhart test of
Section~\ref{sec:operationalizing} predicts that B's alignment
guarantees are tighter than A's by
$\mathrm{Lip}(g)$ times the difference in $\epsnonres$ ceilings
between the two scenarios
on each Lipschitz property $g$.  The cell-wise attribution
tells the deployer \emph{which} infrastructure investment
produced the tightening.  In a transition from A to B, attribute
the gap reduction across:
$\delta_{\mathrm{dormant}}$ reduction (Channel~1 working),
attestation-limited $\delta_{\mathrm{covered}}$ reduction
(Channel~2 working), and any residual $\delta_{\mathrm{covered}}$
proxy-failure-candidate reduction (mixed Channel~1/2 effect or
genuine proxy improvement).  The attribution is what turns a
qualitative claim (``B is better'') into a structurally
precise one (``B has 70\% less $\delta_{\mathrm{dormant}}$ and
60\% less attestation-limited $\delta_{\mathrm{covered}}$, so
the tightening is mostly Channel~1 with a meaningful Channel~2
contribution'').

\subsection{The fungibility collapse: standard economics as macro limit (no channel)}
\label{sec:wex_collapse}

\paragraph{Scenario.}
A representative-agent macroeconomic application: aggregate
consumer welfare, GDP-weighted, on a population of agents
trading approximately fungible goods through liquid markets at
representative-agent scales of aggregation.  The standard
money-maximization model has been empirically successful for
this regime: consumption smoothing under liquidity constraints,
capital allocation across sectors, and intertemporal substitution
at aggregate level all fit reasonably well.  Why?

\paragraph{Computation through the instruments.}
The fungibility-collapse correspondence
(Theorem~\ref{thm:fungibility_collapse}) explains the
phenomenon.  At macro scale on $\Pstd$, the cooperative term
in M6 vanishes, individuation collapses to representative
agents, and the poset degenerates to the wealth axis.  GFM's
$\volL$ on $\Pstd$ reduces to total wealth, and
$\volL$-maximization reduces to expected-wealth maximization.
Consequently, on the observed-and-non-residual subset, the
proxy-truth gap $\epsnonres$ approaches zero: the proxy is
effectively exact when the poset is fungible, because the
five structural features GFM adds beyond $\Pstd$
(non-fungibility, cooperative capabilities, anti-monopolar
reasoning, individuation, residual recognition) all become
operationally irrelevant.

What does \emph{not} go to zero under collapse: $\epsfloor$.
The residual class $\ResS$ remains non-empty.  Household
production, volunteer activity, gift-economy capabilities,
informal economy activity, dark-market activity all remain
outside the channel's reach even when the observed subset's
poset is fully fungible.  Standard economics knows this; SSF
2009 documents the gap at length~\citep{stiglitz2009report}.
The gap is a channel-reach limit, not a poset-structure limit,
and the fungibility-collapse correspondence does not extend
the channel's reach.

\paragraph{Operative channels and instruments.}
\begin{itemize}
\item \textbf{No within-channel channels operative.}  The
  collapse happens to be the regime where Channels~1, 2, 4
  have nothing to tighten because the observed-and-non-residual
  subset's $\epsnonres$ is already approximately zero.
\item \textbf{Channel~3} (individuation) is the only channel
  with content: re-individuating capabilities currently in
  $\ResS$ (treating household production as a distinct
  imputed-time-priced capability, etc.) reduces $\epsfloor$ by
  $\Delta_3^{\mathrm{floor}}$ at the cost of adding
  $\Delta_3^{\mathrm{new}}$ to the within-channel ceiling, with
  net effect $\Delta_3^{\lambda} = \lambda \cdot
  \Delta_3^{\mathrm{floor}} - \Delta_3^{\mathrm{new}}$ on the
  alignment slack.  Channel~3 tightens iff $\Delta_3^{\lambda}
  \geq 0$, i.e., when the imputation's measurement slack does
  not exceed $\lambda$ times the floor reduction.  This is the
  welfare-economics imputation programme
  \citep{stiglitz2009report,aguiar2007measuring} read through
  the GFM lens.
\item \textbf{Instruments:}
  $\delta_{\mathrm{residual}}$ (cell mass at the floor),
  $\epsfloor$ (deployment-specific share).
\end{itemize}

\paragraph{Lesson.}
Standard economics' empirical success at macro scale is
exactly the regime where (i)~the observed poset is
approximately fungible, (ii)~the residual class is small
enough to neglect under representative-agent aggregation, and
(iii)~cooperative / individuation structure is operationally
insignificant.  When any of the three conditions
fails (micro-scale bundle decisions, non-market sectors,
cross-cultural welfare comparison, multi-substrate AI
capability landscapes), the standard model breaks down and the
five structural features become load-bearing
(Section~\ref{sec:fc_features}).  The fungibility-collapse
correspondence makes the breakdown structurally visible: it
tells the reader where the standard model is right (and why)
and where it fails (and through which feature).

The example also illustrates that ``better economic models
$\to$ tighter alignment'' is not a uniform claim across
deployment regimes.  At macro scale on $\Pstd$, alignment is
already tight by Theorem~\ref{thm:fungibility_collapse} on
the observed subset and bounded below by $\lambda \cdot
\epsfloor$ on the residual; further within-channel tightening
has no content because the within-channel gap is already
approximately zero.  The binding constraint is the floor, and
the operational lever is Channel~3 (re-individuation of
non-market capabilities).  This is the same operational lever
welfare economics identified in 2009.

\subsection{Residential relocation: mixed-polarity bundle-for-bundle trades (Channels~3 and~4)}
\label{sec:wex_relocation}

\paragraph{Scenario.}
An agent moves from home~A to home~B at the same nominal
purchase price.  The trade surrenders bundle
$X = \{\mathrm{shelter}_A, \mathrm{neighborhood}_A,
\mathrm{commute}_A, \mathrm{school}_A, \dots\}$ and acquires
bundle $Y = \{\mathrm{shelter}_B, \mathrm{neighborhood}_B,
\mathrm{commute}_B, \mathrm{school}_B, \dots\}$.  Two agents
perform nominally identical priced relocation trades.
Agent~$A$'s pre-trade state has shelter at sub-sufficiency
(foreclosure-driven relocation, with retention degrading
sufficiency and no non-degrading third option available);
agent~$J$'s pre-trade state has shelter at above-sufficiency
(job-driven relocation, with retention non-degrading and
multiple non-degrading alternatives available).

\paragraph{Utility-theory account.}
$\Delta U = U(B) - U(A) > 0$ for both agents (each chose~B).
No further structural decomposition.  Two agents performing
the same priced trade are treated identically under linear
utility; under heterogeneous concave $u_i$, the variance is
attributed to $u_i$'s shape rather than to the agents'
pre-trade need-satisfaction states.  The framework cannot
generate the bifurcation prediction.

\paragraph{Computation through the instruments.}
By
\citep[Definition~3]{lasser2026paper8b}, the polarity
boundary on the shelter dimension distinguishes need-side
(below sufficiency, retention degrading) from want-side
(above sufficiency, retention non-degrading).  Agent~$A$'s
shelter trade is below the polarity boundary on shelter; the
sacrifice magnitude on the shelter component reveals
cost-of-access, captured by NAC
(\citep[Definition~9]{lasser2026paper8b}).  Agent~$J$'s
shelter trade is above the polarity boundary; the sacrifice
magnitude reveals preference, captured by the standard
revealed-sacrifice channel.  The non-shelter components
(neighborhood, commute, school) are above sufficiency for
both agents and contribute want-side sacrifice.

The bundle decomposition gives the per-component split.
Agent~$A$'s split has high need-share:
$\mathrm{NAC}_A(Y) / \mathrm{Sacrifice}_A(X) \gg
\mathrm{NAC}_J(Y) / \mathrm{Sacrifice}_J(X)$.  The
S1-attestation layer's H1 (retention degrading) and H2 (no
non-degrading alternative) classifications place agent~$A$
in the below-sufficiency-with-no-alternative regime
(\citep[Definition~6]{lasser2026paper8b});
agent~$J$ remains in the above-sufficiency regime.  The
classification is what bifurcates the population into the
two regimes Prediction~\ref{pred:mixed_polarity} predicts.

\paragraph{Operative channels and instruments.}
\begin{itemize}
\item \textbf{Channel~3} (individuation): distinguishing
  need-share from want-share on the shelter dimension
  requires individuating ``shelter at or below sufficiency''
  as a capability with NAC accounting separate from
  ``shelter above sufficiency.''  The polarity boundary is
  the individuation cut.
\item \textbf{Channel~4} (bundle decomposition): the per-component
  $\alpha_{n,c}$ decomposition is what enables the
  shelter-vs-non-shelter split.  Without it, the trade is
  observed as a single priced bundle and the bifurcation is
  invisible.
\item \textbf{Instruments:} NAC, S1-attestation, polarity
  boundary, per-event $\alpha_{n,c}$.
\end{itemize}

\paragraph{What the bifurcation enables.}
The framework's structural ability to distinguish agent~$A$
from agent~$J$ on nominally identical priced trades is
operationally significant in several ways, none of which
require taking a normative position on which agent should be
favored:

\begin{enumerate}
\item \emph{Distributional measurement.}  The framework yields
  a population-level decomposition of relocation trades by
  pre-trade need-satisfaction state.  Utility theory cannot
  extract this signal from priced-trade data alone: it sees
  only the price and the chosen direction.  GFM sees the
  fraction of aggregate sacrifice that is need-driven (NAC)
  versus want-driven (above-sufficiency).  This distribution
  is a measurement artifact in its own right, prior to any
  policy conclusion drawn from it.
\item \emph{Disparate-incidence diagnosis.}  Any policy that
  operates on a single proxy of relocation activity
  (relocation tax, mobility subsidy, displacement
  compensation) treats both regimes identically.  The
  bifurcation makes the differential incidence of the
  uniform-proxy policy structurally visible: the policy's
  effect on the high-NAC-share cluster differs from its effect
  on the low-NAC-share cluster, by an amount the framework
  can quantify.  Whether the differential is unjust, intended,
  or acceptable is a normative question the framework does
  not answer; whether it exists is now a question with a
  formal answer.
\item \emph{Structural needs-test.}  The H1+H2 classification
  (retention degrading + no non-degrading alternative) is a
  structural needs-test that operates on the trade itself
  rather than on ad hoc demographic eligibility cutoffs.
  Eligibility for need-targeted assistance, if a deployment
  chooses to target need, follows from the structure of each
  trade rather than from income thresholds, age brackets, or
  household-composition rules.  This does not say the
  framework's needs-test is the right one to use; it says the
  framework provides a needs-test that is structural rather
  than categorical, and a deployment can compare it against
  whatever categorical test it currently uses.
\item \emph{Welfare-aggregation hygiene.}  Aggregate welfare
  measures pooling need-driven and want-driven trades inherit
  a polarity-mixed signal in which sacrifice
  on the high-NAC side reveals cost-of-access while sacrifice
  on the low-NAC side reveals preference.  The bifurcation
  tells the analyst what fraction of the aggregate is
  preference-revealing versus cost-of-access, which is
  information the analyst needs to interpret the aggregate
  honestly.
\end{enumerate}

\paragraph{Tests for the bifurcation.}
Three concrete tests of the bifurcation hypothesis itself:
\begin{enumerate}
\item \emph{Bifurcation test.}  Cross-sectional analysis of
  agents performing similar priced relocations, stratified
  by pre-trade need-attestation.  Prediction: NAC-share
  versus above-sufficiency-share splits cluster bimodally
  on shelter sufficiency rather than continuously on the
  priced trade magnitude.
\item \emph{S1-attestation test.}  H1+H2 classification
  should partition the agent population on shelter
  sufficiency consistently with the bifurcation in test~1;
  agents whose H1+H2 status indicates degrading retention
  with no non-degrading alternative should be the
  high-NAC-share cluster.
\item \emph{Population HHI test.}  The wireheading-consistent
  HHI on relocation sub-categories should distinguish
  forced-relocation regions (high HHI, mass concentrated
  in foreclosure / displacement categories) from
  voluntary-mobility regions (lower HHI, mass distributed
  across job-driven, family-driven, lifestyle-driven
  categories).  This is Prediction~\ref{pred:anti_monopolar}
  applied to the relocation sub-domain.
\end{enumerate}

\paragraph{Goodhart deployment misalignment.}
A deployment that optimizes a coarser proxy than the
formalization warrants (e.g., a relocation-tax framework
optimizing ``average price paid for relocation'', a $\volL$
proxy with no need-share / want-share distinction) treats
both regimes identically.  Under Theorem~\ref{thm:goodhart},
the proxy-truth gap on this deployment includes the entire
below-sufficiency contribution to NAC, and the misalignment
magnitude is bounded by $\mathrm{Lip}(g) \cdot \epsnonres$ on
the relocation-tax outcome.  Sharper economics (admitting
bundle decomposition and polarity boundary into the framework)
delivers a tighter alignment to the true distribution of
welfare consequences.  We emphasize: this does not say the
deployment \emph{should} treat the two regimes differently;
it says that if the deployment treats them identically, the
cost in alignment-target accuracy is bounded by
Theorem~\ref{thm:goodhart} and the bound is not zero.

A reading caveat on the bound.  The misalignment magnitude
$\mathrm{Lip}(g) \cdot \epsnonres$ is a function of the
deployment's chosen individuation: a deployment that has
\emph{already} performed the shelter-at-vs-above-sufficiency
individuation (drawn the polarity boundary) has a smaller
$\epsnonres$ that distinguishes the two regimes; a deployment
that has \emph{not yet} performed the individuation has a
larger $\epsnonres$ on a coarser capability space that does not
see the bifurcation at all, with a correspondingly larger but
structurally different bound.  The example partly motivates
the individuation move: the bound's bite depends on the
individuation having occurred, but the comparison between
``with bifurcation'' and ``without bifurcation'' bounds is what
shows the deployer the cost of staying on the coarser
individuation.

This is the canonical case where Channels~3 and~4 jointly
deliver alignment-tightening that no single channel could
deliver.  Channel~3 alone (individuating shelter sufficiency)
would not separate the trades' shelter contributions from
their non-shelter contributions; Channel~4 alone (bundle
decomposition without polarity discipline) would split the
trade into components but not classify the shelter
component's polarity.  Per the sequential-intervention regime
of Remark~\ref{rem:ch34_interaction}, we treat the example as
the deployment first fixing the individuation discipline that
distinguishes shelter sufficiency (Channel~3, contributing its
$\lambda$-weighted net $\Delta_3^{\lambda}$), then performing
bundle decomposition on the resulting individuation
(Channel~4, contributing $\Delta_4$).  Under the sequential
regime, Proposition~\ref{prop:channel_composition} applies
directly; co-evolution between the two channels would require
the correction terms flagged in
Remark~\ref{rem:ch34_interaction}.

\subsection{What the four examples demonstrate}
\label{sec:wex_summary}

The four examples cover four scales (micro, AI alignment,
macro, bundle-for-bundle) and exercise four combinations of
the channels:

\begin{table}[ht]
\centering
\footnotesize
\setlength{\tabcolsep}{4pt}
\begin{tabular}{>{\raggedright\arraybackslash}p{0.20\linewidth}
                >{\raggedright\arraybackslash}p{0.20\linewidth}
                >{\raggedright\arraybackslash}p{0.30\linewidth}
                >{\raggedright\arraybackslash}p{0.22\linewidth}}
\toprule
\textbf{Example} & \textbf{Scale} & \textbf{Operative channels} &
\textbf{Headline prediction} \\
\midrule
\S\ref{sec:wex_boat}: Boat &
Micro &
Ch.~4 (bundle decomp) + Ch.~2 (attestation) &
Complements amplify cooperative excess; cross-agent variance \\
\addlinespace
\S\ref{sec:wex_goodhart_gradient}: Goodhart gradient &
AI alignment &
Ch.~1 (observation density) + Ch.~2 (attestation) &
$\sim$8$\times$ ceiling tightening from A to B \\
\addlinespace
\S\ref{sec:wex_collapse}: Fungibility collapse &
Macro &
None within-channel; Ch.~3 (individuation) for floor &
$\epsnonres \to 0$ but $\epsfloor$ binds \\
\addlinespace
\S\ref{sec:wex_relocation}: Residential relocation &
Bundle-for-bundle &
Ch.~3 (individuation) + Ch.~4 (bundle decomp) &
Bifurcation by pre-trade sufficiency \\
\bottomrule
\end{tabular}
\caption{The four worked examples and their channel composition.
The examples are deliberately chosen to span the structural
variety: a single-event prediction (boat), a cross-domain
infrastructure comparison (Goodhart gradient), a regime-limit
case (fungibility collapse), and a deployment-misalignment case
(residential relocation).  Together they exercise all six
instruments of Table~\ref{tab:instruments} and demonstrate that
the channel composition of
Proposition~\ref{prop:channel_composition} is operationally
visible in concrete data structures.}
\label{tab:wex_summary}
\end{table}

The examples support three reading-level claims about the
formalization.  First, the channels of
Section~\ref{sec:lineage} are operationally distinct: each
example's binding channels are different combinations, and the
examples exercise the full channel space across the four cases.
Second, the gap decomposition of
\citep[Need~Sufficiency]{lasser2026paper8b} plays its predicted
selector role: in each example, the cell-wise attribution tells
the deployer which channel intervention applies.  Third, the
Goodhart machinery of Section~\ref{sec:goodhart} is testable in
practice: Section~\ref{sec:wex_goodhart_gradient}'s cross-domain
comparison is the canonical structural test, and
Section~\ref{sec:wex_relocation}'s deployment-misalignment is
the canonical concrete consequence.

The examples do not constitute empirical validation.  They are
\emph{methodology demonstrations}: they show how the formalization
applies, what it predicts, and what would falsify it.  Empirical
validation of the predictions across actual deployments is
deferred work; the contribution at this stage is to identify the
predictions clearly enough that they can be tested and to
establish the methodology template through which the testing
operates.
